\documentclass[11pt,letterpaper]{article}
\usepackage[T1]{fontenc}
\usepackage[utf8]{inputenc}
\usepackage{newpxtext}
\usepackage[margin=1in,headheight=15pt]{geometry}
\usepackage{amsmath,amsthm,mathtools}
\usepackage{newpxmath}
\usepackage{microtype,booktabs,array,longtable,enumitem}
\usepackage[dvipsnames]{xcolor}
\usepackage{fancyhdr}
\usepackage[colorlinks=true,linkcolor=MidnightBlue,citecolor=MidnightBlue,urlcolor=MidnightBlue,breaklinks=true]{hyperref}
\usepackage[nameinlink,noabbrev]{cleveref}
\hypersetup{pdftitle={A Corona Obstruction and Infinite-Dimensional Spectral Fibres in Surcomplex Analysis},pdfauthor={Research report prepared for Vladimir Reshetnikov},pdfsubject={Hahn holomorphic algebras, Bezout identities, ultrafilter primes and geometric local rings}}
\pagestyle{fancy}\fancyhf{}
\fancyhead[L]{\small Surcomplex corona and spectral fibres}
\fancyhead[R]{\small Research article}
\fancyfoot[C]{\thepage}
\renewcommand{\headrulewidth}{0.3pt}
\setlength{\parskip}{0.28em}
\setlength{\parindent}{1.2em}
\setlength{\emergencystretch}{2.5em}
\setlist{nosep,leftmargin=2em}
\allowdisplaybreaks[2]
\numberwithin{equation}{section}
\newtheorem{theorem}{Theorem}[section]
\newtheorem{proposition}[theorem]{Proposition}
\newtheorem{lemma}[theorem]{Lemma}
\newtheorem{corollary}[theorem]{Corollary}
\theoremstyle{definition}
\newtheorem{definition}[theorem]{Definition}
\newtheorem{example}[theorem]{Example}
\theoremstyle{remark}
\newtheorem{remark}[theorem]{Remark}
\newcommand{\C}{\mathbb C}
\newcommand{\R}{\mathbb R}
\newcommand{\Q}{\mathbb Q}
\newcommand{\Z}{\mathbb Z}
\newcommand{\N}{\mathbb N}
\newcommand{\No}{\mathbf{No}}
\newcommand{\SC}{\mathbf{No}[i]}
\newcommand{\OO}{\mathcal O}
\newcommand{\HH}{\mathcal H}
\newcommand{\BB}{\mathcal B}
\newcommand{\UU}{\mathcal U}
\newcommand{\mm}{\mathfrak m}
\newcommand{\pp}{\mathfrak p}
\newcommand{\vH}{v_{\!H}}
\newcommand{\Dsharp}{D^{\sharp}_{\Gamma}}
\newcommand{\ev}{\operatorname{ev}}
\newcommand{\supp}{\operatorname{supp}}
\newcommand{\Spec}{\operatorname{Spec}}
\newcommand{\Max}{\operatorname{Max}}
\newcommand{\Frac}{\operatorname{Frac}}
\newcommand{\st}{\operatorname{st}}
\newcommand{\ord}{\operatorname{ord}}
\newcommand{\diag}{\operatorname{diag}}
\newcommand{\rk}{\operatorname{rank}}
\newcommand{\norm}[1]{\left\lVert#1\right\rVert}
\newcommand{\abs}[1]{\left\lvert#1\right\rvert}
\newcolumntype{L}[1]{>{\raggedright\arraybackslash}p{#1}}

\begin{document}
\begin{titlepage}
\centering
\vspace*{0.25in}
{\large\scshape Surreal and surcomplex research\par}
\vspace{0.32in}
{\Huge\bfseries A Corona Obstruction\par}
\vspace{0.10in}
{\LARGE\bfseries and Infinite-Dimensional Spectral Fibres\par}
\vspace{0.12in}
{\LARGE in Surcomplex Analysis\par}
\vspace{0.25in}
{\large Global Hahn support, exact B\'ezout criteria,\par and regular local rings at every geometric point\par}
\vspace{0.25in}
{\large Prepared for Vladimir Reshetnikov\par}
{\normalsize September 21, 2026\par}
\vspace{0.25in}
\begin{minipage}{0.96\textwidth}
\begin{abstract}
\noindent
Fix the ordinary unit disk $D$ and the set-sized surcomplex field
$K=\C((t^{\R}))\subset\No[i]$, with $t^r=\omega^{-r}$. We investigate
$\HH=\OO(D)((t^{\R}))$: the algebra of Hahn series whose coefficients are
holomorphic on the \emph{same fixed disk}. It is a complete uniform
non-Archimedean Banach $K$-algebra, and its norm is the supremum of its
values at actual finite surcomplex points over $D$. Nevertheless, we give
explicit $h,F\in\HH$ with norms at most one and
\[
 \inf_{z\in D^{\sharp}}\max\{|h(z)|_v,|F(z)|_v\}=e^{-2}>0
\]
for which $hU+FV=1$ has no solution in $\HH$; the distance from $1$ to
$(h,F)$ is exactly one. The obstruction survives every enlargement of the
Hahn value group. An exact criterion explains it: modulo a fixed ordinary
simple discrete divisor, a nowhere-zero value family is invertible
precisely when its set of leading valuations is finite.

We also construct a continuum-sized strictly ordered family of prime
ideals of $\HH$, all contracting to a single free maximal ideal of
$\OO(D)$. Thus one spectral fibre is an infinite-dimensional domain; the displayed
primes have the same norm closure and lie below an explicit
coefficient-reduction maximal ideal. They are contained in no geometric
evaluation maximal ideal, whereas
localization at \emph{every} actual surcomplex evaluation point is a
discrete valuation ring with completion $K[[X]]$. The results separate
geometric regularity from the much larger global algebraic spectrum and
correct a fixed-domain structure assertion in the supplied repository.
The explicit corona example and prime-fibre construction have real forms
inside surreal-valued analysis.
\end{abstract}
\end{minipage}
\vfill
\begin{minipage}{0.96\textwidth}\small
\textbf{Status and novelty.} This is a proof-based research draft, not a
claim to solve a named published conjecture. Hahn summability, ordinary
interpolation, and growth-rate constructions of prime ideals have
substantial prior literature. The proposed contribution is the precise
fixed-domain theorem package above, particularly the normed corona
counterexample, its exact finite-valuation criterion, and the
infinite-dimensional fibre over one classical maximal ideal. No matching
statement was located in the targeted search described in
\cref{sec:prior}. Priority is not certified; the proofs have not been
independently refereed or checked in a proof assistant.
\end{minipage}
\end{titlepage}

\pagenumbering{roman}
\tableofcontents
\clearpage
\pagenumbering{arabic}

\section{The question and the principal results}\label{sec:introduction}

\subsection{A global question that pointwise analysis cannot answer}
The algebraic closure $\No[i]$ of the surreal field supplies infinitesimal
and infinite complex scales. One rigorous way to use these scales in
analysis is to fix a Hahn field and let ordinary holomorphic coefficient
functions vary in it. The resulting objects retain ordinary analytic
information and exact infinitesimal information simultaneously. They also
have a compatibility condition that ordinary complex analysis does not:
the exponents used across the whole coefficient domain must form one
well-ordered set.

This condition raises a concrete question. Suppose two such functions
have no common zero at any actual surcomplex point of their domain. Does
there exist a global B\'ezout identity between them? Does a uniform
positive lower bound on their pointwise sizes help? And, if every
geometric local ring is regular of dimension one, does the global
algebraic spectrum have dimension one as well?

For the fixed-domain Hahn algebra studied here, the answer to each of
these proposed implications is negative. The obstruction is not a
singularity of an individual coefficient and not an inability to solve a
finite system. Every finite interpolation problem in our example is
solvable. The obstruction appears when inverse values force a decreasing
infinite sequence of exponents into one global support.

\subsection{The precise ring}
Throughout the article,
\[
 D=\{z\in\C:|z|<1\},\qquad
 K_\Gamma=\C((t^\Gamma)),\qquad
 \HH_\Gamma=\OO(D)((t^\Gamma)),
\]
where $\Gamma$ is a set-sized ordered abelian group. For its surcomplex
interpretation we take $\Gamma\subset(\No,+)$ and interpret
$t^\gamma=\omega^{-\gamma}$. Divisibility of $\Gamma$ is not required for
the main algebraic arguments. The normed and prime-fibre theorems use
$\Gamma=\R$; write $K=K_\R$ and $\HH=\HH_\R$ there.

An element $F=\sum_{\gamma\in S}f_\gamma t^\gamma$ belongs to
$\HH_\Gamma$ exactly when $S$ is well ordered and every $f_\gamma$ is
holomorphic on the \emph{whole} disk $D$. The coefficients need not have
uniform ordinary bounds. This is neither a ring of germs at one ordinary
point nor a Tate algebra. Taylor evaluation is defined on
\[
 D^\sharp_\Gamma=\{c+\varepsilon:c\in D,\ \varepsilon\in K_\Gamma,
                                  \ v(\varepsilon)>0\},
\]
where $v(0)=+\infty$, so ordinary points are included.

\subsection{An explicit divisor and the theorem package}
Set, for $n\ge1$,
\begin{equation}\label{eq:cardinal-data}
 a_n=1-\frac1n,\qquad
 h(z)=\sin\frac{\pi}{1-z},\qquad
 e_n(z)=\frac{h(z)}{\pi n^2(-1)^n(z-a_n)}.
\end{equation}
The removable values are understood. We prove below that $h$ has exactly
the simple zeros $a_n$ in $D$, and that $e_n(a_m)=\delta_{nm}$.

\begin{theorem}[Exact B\'ezout criterion]\label{thm:main-bezout}
For every $F\in\HH_\Gamma$, the following conditions are equivalent:
\begin{enumerate}
\item there are $U,V\in\HH_\Gamma$ with $hU+FV=1$;
\item $F(a_n)\ne0$ for every $n$ and
      $\{v(F(a_n)):n\ge1\}$ is finite.
\end{enumerate}
Solvability of this identity is unchanged by any ordered value-group
extension $\Gamma\subset\Gamma'$.
\end{theorem}

\begin{theorem}[Uniform corona failure]\label{thm:main-corona}
In $\HH=\OO(D)((t^\R))$, put
\begin{equation}\label{eq:F-main}
 \gamma_n=2-\frac1n,\qquad
 F(z)=\sum_{n\ge1}t^{\gamma_n}e_n(z).
\end{equation}
The Hahn algebra has the complete multiplicative norm
$\norm{G}_H=e^{-\vH(G)}$, and
\[
 \norm{G}_H=\sup_{z\in D^\sharp_\R}|G(z)|_v,
 \qquad |x|_v=e^{-v(x)}.
\]
For the pair in \eqref{eq:cardinal-data}--\eqref{eq:F-main},
\begin{align*}
 \norm h_H&=1,&\norm F_H&=e^{-1},\\
 \inf_{z\in D^\sharp_\R}\max\{|h(z)|_v,|F(z)|_v\}&=e^{-2},&
 \inf_{U,V\in\HH}\norm{1-hU-FV}_H&=1.
\end{align*}
In particular $(h,F)$ is a proper, nonprincipal two-generated ideal with
empty geometric zero set. The pair admits local B\'ezout identities on
ordinary neighborhoods of every point of $D$, but no global one.
\end{theorem}

\begin{theorem}[Infinite-dimensional fibres and geometric regularity]
\label{thm:main-spectrum}
Fix a nonprincipal ultrafilter $\UU$ on $\N_{\ge1}$. For each real $r\ge0$
define
\begin{equation}\label{eq:Pr-intro}
 P_r=\left\{G\in\HH:
   \{n:v(G(a_n))>m n^r\}\in\UU\text{ for every }m\ge1\right\}.
\end{equation}
These are proper prime ideals, and $P_s\subsetneq P_r$ whenever $r<s$.
All contract to the same maximal ideal
\[
 \mm_\UU=\{f\in\OO(D):\{n:f(a_n)=0\}\in\UU\}.
\]
Thus the fibre of $\Spec\HH\to\Spec\OO(D)$ over $\mm_\UU$ has infinite
Krull dimension and a chain of cardinality $2^{\aleph_0}$. There is a
maximal ideal $M$ above $\mm_\UU$ for which $\dim\HH_M=\infty$.

In contrast, for every $a\in D^\sharp_\Gamma$, the kernel of evaluation is
$\mm_a=(z-a)$, the localization $(\HH_\Gamma)_{\mm_a}$ is a discrete
valuation ring with residue field $K_\Gamma$, and its completion is
$K_\Gamma[[X]]$. None of the primes $P_r$ is contained in any such
geometric maximal ideal.
\end{theorem}

The proofs do not rely on the repository's unverified structure results.
They use ordinary holomorphic interpolation, the classical Hahn support
lemma, and explicit algebra. The restricted interpolation quotient itself
has a close predecessor in the repository's global-divisor report
\cite{RepoGlobal}; we credit it and prove the simple-divisor case needed
here rather than relabel it as a discovery.

\section{Hahn support and surcomplex evaluation}\label{sec:hahn}

\subsection{Formal summability, not a limit of partial sums}
For a commutative ring $A$, a Hahn series over $\Gamma$ is a function
$\Gamma\to A$ with well-ordered support, written $\sum a_\gamma t^\gamma$.
For a family of such series $(F_i)_{i\in I}$, \emph{strong summability}
means that the union of supports is well ordered and that, for each
$\gamma$, only finitely many $F_i$ have a nonzero coefficient at
$\gamma$. Its sum is then defined coefficientwise.

The distinction between the index $i$ of a summation and a spatial index
$n$ is important. A coefficient sequence $(c_n)\in\C^\N$ is allowed to
have infinitely many nonzero entries. We never add those entries merely
because they occur at the same Hahn exponent.

We use the following classical support theorem, usually associated with
the Hahn--Mal'cev--Neumann construction \cite{Hahn,Neumann}. A recent
account with a different proof of positive-support inversion is
\cite{Poonen}; formal summability in more general algebras is developed in
\cite{BagayokoEtAl}.

\begin{lemma}[Classical Hahn--Neumann support lemma]\label{lem:neumann}
Let $S,T\subset\Gamma$ be well ordered and let
$E\subset\Gamma_{>0}$ be well ordered.
\begin{enumerate}
\item $S+T$ is well ordered, and a fixed element has only finitely many
      representations $s+t$ with $(s,t)\in S\times T$.
\item The monoid $E^*=\{0\}\cup E\cup(E+E)\cup\cdots$ is well ordered.
      Each element has only finitely many representations as a finite
      ordered string of elements of $E$.
\item If $U$ has support contained in $E$, then
      $\sum_{k\ge0}(-U)^k$ is strongly summable and is the inverse of
      $1+U$.
\end{enumerate}
These statements hold over arbitrary ordered abelian groups, not only
Archimedean ones.
\end{lemma}

The first assertion is also easily seen from sequences: an infinite
family of representations of a fixed sum would produce an increasing
subsequence in one summand and a decreasing subsequence in the other.
The positive-monoid assertion is the substantive classical input. We do
not replace it by the false high-rank assertion that $k\delta$ eventually
exceeds every element of $\Gamma$.

\subsection{Valuations on the scalar field and the coefficient algebra}
For nonzero $F\in\HH_\Gamma$, let
\[
 \vH(F)=\min\supp F.
\]
The disk is connected, so $\OO(D)$ is a domain. Consequently
\begin{equation}\label{eq:global-valuation}
 \vH(FG)=\vH(F)+\vH(G),\qquad
 \vH(F+G)\ge\min\{\vH(F),\vH(G)\}.
\end{equation}
The first identity follows because the leading coefficient of the product
is the product of the two nonzero leading coefficient functions. Thus
$\HH_\Gamma$ is a domain. The scalar Hahn field $K_\Gamma$ embeds as the
subalgebra of constant coefficient functions.

\begin{lemma}[A sufficient global unit criterion]\label{lem:global-unit}
If the leading coefficient $f_\alpha$ of $F\in\HH_\Gamma$ is nowhere
zero on $D$, then $F$ is a unit in $\HH_\Gamma$.
\end{lemma}
\begin{proof}
Write $F=t^\alpha f_\alpha(1+U)$. The coefficient functions of $U$ are
holomorphic on $D$, since $1/f_\alpha$ is, and the support of $U$ is
strictly positive. Apply \cref{lem:neumann} and multiply the resulting
inverse by $t^{-\alpha}/f_\alpha$.
\end{proof}

\subsection{Evaluation at actual infinitesimal displacements}
Let $a=c+\varepsilon\in D^\sharp_\Gamma$. Define
\begin{equation}\label{eq:evaluation}
 F(a)=\sum_{\gamma\in\supp F}\sum_{k\ge0}
     \frac{f_\gamma^{(k)}(c)}{k!}\,t^\gamma\varepsilon^k.
\end{equation}
When $\varepsilon=0$, the inner sum has just its constant term. Otherwise
let $E=\supp\varepsilon\subset\Gamma_{>0}$. The terms in
\eqref{eq:evaluation} have support in $\supp F+E^*$, and the finiteness
assertions in \cref{lem:neumann} justify every coefficient. Thus no common
ordinary Taylor radius for all $f_\gamma$ is needed at this step.

\begin{proposition}[Evaluation and its basic estimate]\label{prop:eval}
The map $\ev_a:\HH_\Gamma\to K_\Gamma$ is a surjective $K_\Gamma$-algebra
homomorphism. For all $F$,
\begin{equation}\label{eq:eval-bound}
 v(F(a))\ge\vH(F).
\end{equation}
For an ordinary coefficient function $f$ one has $v(f(a))\ge0$, and if
$f(c)\ne0$, then $v(f(a))=0$.
\end{proposition}
\begin{proof}
Taylor's product formula and the finite rearrangements licensed by
\cref{lem:neumann} show that \eqref{eq:evaluation} preserves products and
sums. It fixes the scalar Hahn field, so it is surjective. Each
nonzero contribution to \eqref{eq:evaluation} has exponent at least
$\vH(F)$, which proves the estimate. When $f(c)\ne0$, the exponent-zero
term is nonzero and all remaining contributions have positive exponent.
\end{proof}

This is a Taylor extension of ordinary functions to finite surcomplex
points. It is not the Berarducci--Mantova derivation, and no sine of an
infinite surreal argument is being defined. In particular, the formula
for $h$ means an ordinary holomorphic coefficient function followed by
\eqref{eq:evaluation}.

The embedding into $\No[i]$ is the familiar normal-form interpretation of
Hahn series. General background on surreal fields and their normal forms
is given in \cite{Gonshor,vdDE}. All rings, supports, function domains, and
spectra used in the proofs are sets. No spectrum of a proper-class field
is required.

\section{The fixed-domain algebra is a uniform Banach algebra}
\label{sec:banach}
For this section $\Gamma=\R$. Give $K$ its intrinsic absolute value
$|x|_v=e^{-v(x)}$ and set $\norm F_H=e^{-\vH(F)}$, with
$\norm0_H=0$. This is \emph{not} the ordered surreal modulus of a
surcomplex number. In particular every nonzero ordinary complex scalar
has valuation norm one.

\begin{proposition}[Completeness and exact pointwise norm]
\label{prop:banach}
$\HH$ is a complete normed $K$-algebra. Its norm is multiplicative, hence
power-multiplicative, and
\begin{equation}\label{eq:supnorm}
 \norm F_H=\sup_{a\in D^\sharp_\R}|F(a)|_v.
\end{equation}
For nonzero $F$, the supremum is attained at an ordinary point of $D$.
\end{proposition}
\begin{proof}
The norm axioms and multiplicativity follow from
\eqref{eq:global-valuation}; scalar homogeneity follows by regarding
$K$ as constant functions. To prove completeness, start with a Cauchy
sequence and select a subsequence $(G_j)$ such that
\[
 \vH(G_{j+1}-G_j)>j\qquad(j\ge1).
\]
Put $\Delta_j=G_{j+1}-G_j$. The family consisting of $G_1$ and all
$\Delta_j$ is strongly summable. Indeed, below any fixed real bound
$B$, only finitely many supports of the $\Delta_j$ can occur. A descending
sequence in their union would be bounded above by its first term, and
would therefore lie in a finite union of well-ordered supports, an
impossibility. At each exponent only finitely many summands contribute.
The coefficient functions of the sum are finite sums of functions in
$\OO(D)$, so its sum $G$ lies in $\HH$. The tails have valuation tending
to $+\infty$, hence $G_j\to G$. The original Cauchy sequence has the same
limit.

The inequality $\sup|F(a)|_v\le\norm F_H$ is
\eqref{eq:eval-bound}. If $\alpha=\vH(F)$, choose an ordinary $c\in D$
where $f_\alpha(c)\ne0$. Then $v(F(c))=\alpha$, proving equality and
attainment.
\end{proof}

Thus $\HH$ is a uniform non-Archimedean Banach algebra in the literal
normed-algebra sense. It is not an affinoid algebra and is not the
classical algebra $H^\infty(D)$ of ordinary bounded holomorphic functions.
Our corona counterexample will occur \emph{inside this complete algebra};
incompleteness is not its explanation.

\begin{lemma}[Unit neighborhood and a distance obstruction]
\label{lem:distance}
If $\norm E_H<1$, then $1-E$ is a unit in $\HH$. Consequently every proper
ideal $I\subset\HH$ satisfies
\begin{equation}\label{eq:distance}
 \inf_{A\in I}\norm{1-A}_H=1.
\end{equation}
In particular, the closure of a proper ideal remains proper, and every
maximal ideal is closed.
\end{lemma}
\begin{proof}
Positive support gives the inverse $\sum_{j\ge0}E^j$ by
\cref{lem:neumann}. Alternatively this geometric series converges in the
norm since $\norm E_H<1$. If $A\in I$ and $\norm{1-A}_H<1$, then $A$ is
a unit, contradicting properness. The opposite inequality in
\eqref{eq:distance} follows by choosing $A=0$. The closure assertions
follow immediately.
\end{proof}

\begin{remark}[Two different infinite sums]
The completeness proof uses differences whose valuations tend to
$+\infty$. The series $F$ in \eqref{eq:F-main} instead has exponents
increasing to $2$. It is a valid Hahn series, but its finite partial sums
do not converge to it in $\norm{\cdot}_H$: the remainder after $N$ terms
has norm $e^{-\gamma_{N+1}}\to e^{-2}$, not zero. Neither assertion
concerns the fine topology on the proper class $\No[i]$.
\end{remark}

\section{A discrete divisor and the exact interpolation quotient}
\label{sec:interpolation}

\subsection{The cardinal functions}
\begin{lemma}\label{lem:cardinal}
The zeros of $h$ in \eqref{eq:cardinal-data} are precisely the $a_n$,
each with multiplicity one. Moreover
\[
 h'(a_n)=\pi n^2(-1)^n,\qquad e_n(a_m)=\delta_{nm}.
\]
\end{lemma}
\begin{proof}
For $w=1/(1-z)$, the inequality $|z|<1$ is equivalent to
$|w-1|<|w|$, hence to $\Re w>1/2$. The zeros of $\sin(\pi w)$ in that
half-plane are exactly $w=n$ for positive integers $n$. The derivative
is
\[
 h'(z)=\frac{\pi}{(1-z)^2}\cos\frac{\pi}{1-z},
\]
which gives the stated nonzero values. Removing the simple zero in the
quotient defining $e_n$ gives value one at $a_n$ and zero at the other
$a_m$.
\end{proof}

\subsection{An ordinary interpolation proof on the same disk}
The following is classical; a direct proof is included to make the
common-domain requirement transparent. Ordinary discrete interpolation
and its relation to free ideals appear in \cite{Henriksen52,Henriksen53}.

\begin{lemma}[Unrestricted complex interpolation]\label{lem:ordinary-interp}
For every sequence $(c_n)\in\C^\N$ there exists $f\in\OO(D)$ such that
$f(a_n)=c_n$ for all $n$.
\end{lemma}
\begin{proof}
For $n\ge3$, set $r_n=1-2/n<a_n$. Choose a nonnegative integer $N_n$
large enough that
\[
 \sup_{|z|\le r_n}
 \left|c_ne_n(z)\left(\frac{1+z}{1+a_n}\right)^{N_n}\right|
 \le2^{-n}.
\]
This is possible because $|1+z|/(1+a_n)\le(1+r_n)/(1+a_n)<1$ on the
closed disk in question. Take any nonnegative $N_1,N_2$. The series
\[
 f(z)=\sum_{n\ge1}c_ne_n(z)
                    \left(\frac{1+z}{1+a_n}\right)^{N_n}
\]
converges locally uniformly on $D$, since $r_n\to1$. It defines a
holomorphic function there. At $a_m$, the only nonzero term is the
$m$th, whose extra factor is one. Hence $f(a_m)=c_m$.
\end{proof}

This proof permits arbitrarily large prescribed coefficients. It does
not claim ordinary uniform boundedness on $D$, and it uses ordinary
complex convergence only to construct individual coefficient functions.

\subsection{The common-support product}
\begin{definition}\label{def:B}
Define
\begin{equation}\label{eq:Bdef}
 \BB_\Gamma=\left\{(x_n)\in K_\Gamma^\N:
                 \bigcup_{n\ge1}\supp x_n\text{ is well ordered}\right\}.
\end{equation}
It carries coordinatewise addition and multiplication.
\end{definition}
The Hahn support lemma makes $\BB_\Gamma$ a ring. Equivalently,
\[
 \BB_\Gamma=(\C^\N)((t^\Gamma)),
\]
where $\C^\N$ is the ordinary product ring. This equivalence keeps the
spatial coordinate index separate from summation indices.

\begin{theorem}[Fixed-divisor interpolation quotient]\label{thm:quotient}
The map
\[
 q:\HH_\Gamma\longrightarrow\BB_\Gamma,\qquad
 q(F)=(F(a_n))_{n\ge1}
\]
is a surjective $K_\Gamma$-algebra homomorphism, and
\begin{equation}\label{eq:quotient}
 \ker q=h\HH_\Gamma,\qquad
 \HH_\Gamma/(h)\simeq\BB_\Gamma.
\end{equation}
Every value family has a lift with support contained in the union of its
coordinate supports.
\end{theorem}
\begin{proof}
At an ordinary point evaluation merely evaluates the coefficient
functions. Thus $\supp F(a_n)\subseteq\supp F$, and $q$ takes values in
$\BB_\Gamma$.

Given $x=(x_n)\in\BB_\Gamma$, let $S=\bigcup_n\supp x_n$ and write
$x_n=\sum_{\gamma\in S}c_{\gamma,n}t^\gamma$. For each $\gamma\in S$
apply \cref{lem:ordinary-interp} to $(c_{\gamma,n})_n$, choosing
$f_\gamma\in\OO(D)$ with these values. Then
$F=\sum_{\gamma\in S}f_\gamma t^\gamma$ is a lift. The set of choices is
set-sized. No sum of different coefficient functions at one exponent is
being taken here.

If $q(F)=0$, every coefficient $f_\gamma$ vanishes at every $a_n$.
Since $h$ has exactly these simple zeros, $f_\gamma/h$ extends
holomorphically throughout $D$. Dividing coefficientwise gives
$F/h\in\HH_\Gamma$ with the same support. This proves the kernel
assertion; the reverse inclusion is immediate.
\end{proof}

This is the simple, undeformed case of the support-restricted Chinese
remainder construction described in the repository's global-divisor
report \cite[\S10]{RepoGlobal}. The quotient is a restricted product,
not all of $K_\Gamma^\N$. The new conclusions below come from calculating
its units exactly and studying valuation growth in its prime ideals.

\begin{proposition}[The quotient norm]\label{prop:Bnorm}
For $\Gamma=\R$, $\BB=\BB_\R$ is a closed isometric unital subalgebra of
$\ell^\infty(\N,K)$ with norm
\[
 \norm{x}_{\BB}=\sup_n|x_n|_v
  =\exp\left(-\min\bigcup_n\supp x_n\right)\quad(x\ne0).
\]
The isomorphism \eqref{eq:quotient} identifies this norm with the quotient
norm from $\HH/(h)$, and $h\HH$ is closed.
\end{proposition}
\begin{proof}
A common well-ordered support has a minimum, giving boundedness of the
coordinate norms and the displayed formula. Evaluation is contractive,
so every lift has norm at least $\norm x_\BB$. The support-preserving
lift in \cref{thm:quotient} attains that lower bound. Multiplication by
$h$ is an isometry of the complete space $\HH$, since $\vH(h)=0$.
Its image is therefore closed, and the quotient is complete. This also
makes the isometric image of $\BB$ in $\ell^\infty(\N,K)$ closed.
\end{proof}

\section{Finite valuation spectrum: units, matrices, and B\'ezout}
\label{sec:units}

\subsection{An order-theoretic necessity}
\begin{lemma}\label{lem:two-orders}
A subset $V$ of a linearly ordered set is finite if both its inherited
order and its reverse order are well orders.
\end{lemma}
\begin{proof}
If $V$ is infinite and well ordered, repeatedly taking the least element
not yet selected constructs an infinite strictly increasing sequence.
That sequence is strictly decreasing for the reverse order, contradicting
its being a well order.
\end{proof}

\begin{theorem}[Exact unit criterion in the common-support product]
\label{thm:units}
Let $x=(x_n)\in\BB_\Gamma$. Then
\begin{equation}\label{eq:unit-criterion}
 x\in\BB_\Gamma^\times
 \quad\Longleftrightarrow\quad
 x_n\ne0\ \text{for every }n
 \quad\text{and}\quad
 V(x)=\{v(x_n):n\ge1\}\text{ is finite}.
\end{equation}
\end{theorem}
\begin{proof}
Necessity is more restrictive than coordinatewise invertibility. If $x$
has an inverse $y\in\BB_\Gamma$, then $y_n=x_n^{-1}$, and in particular
all $x_n$ are nonzero. Put $S=\bigcup_n\supp x_n$ and
$T=\bigcup_n\supp y_n$. The set $V(x)$ is a subset of $S$, so it is well
ordered. The set $-V(x)$ is a subset of $T$, so it is also well ordered.
Apply \cref{lem:two-orders} to deduce that $V(x)$ is finite.

Conversely, suppose all coordinates are nonzero and $V=V(x)$ is finite.
For each $n$ write
\[
 x_n=c_nt^{\nu_n}(1+u_n),\qquad c_n\in\C^\times,\quad \nu_n\in V.
\]
The supports of all $u_n$ lie in
\begin{equation}\label{eq:Efinite}
 E=\bigcup_{\nu\in V}\bigl((S-\nu)\cap\Gamma_{>0}\bigr).
\end{equation}
A finite union of well-ordered subsets of a linearly ordered set is well
ordered, so $E$ is well ordered and positive. By \cref{lem:neumann},
\[
 x_n^{-1}=c_n^{-1}t^{-\nu_n}\sum_{k\ge0}(-u_n)^k
\]
is defined, and all these inverse supports lie in
\begin{equation}\label{eq:inverse-support}
 \bigcup_{\nu\in V}(-\nu+E^*).
\end{equation}
This finite union is well ordered. Hence $(x_n^{-1})_n\in\BB_\Gamma$,
which proves sufficiency and supplies an explicit support certificate.
\end{proof}

The nontrivial point is \emph{finiteness}, not boundedness, of the leading
valuation set. Over $\Gamma=\R$, an infinite increasing bounded sequence
of leading valuations violates the criterion just as decisively as an
unbounded sequence does.

\begin{corollary}[Failure of inverse-closedness]\label{cor:inverseclosed}
The closed subalgebra $\BB_\R\subset\ell^\infty(\N,K)$ is not
inverse-closed. Specifically,
\[
 x_n=t^{2-1/n}
\]
defines an element of $\BB_\R$ whose coordinatewise inverse is bounded
in $\ell^\infty(\N,K)$ but does not belong to $\BB_\R$.
\end{corollary}
\begin{proof}
The common support $\{2-1/n:n\ge1\}$ is well ordered, and the inverse
coordinate norms are $e^{2-1/n}<e^2$. Their required support is
$\{-2+1/n:n\ge1\}$, a strictly decreasing infinite sequence. Apply
\cref{thm:units}.
\end{proof}

\subsection{A finite matrix version}
\begin{corollary}[Matrix inverse criterion]\label{cor:matrix}
Let $A\in M_d(\BB_\Gamma)$, and let $A_n\in M_d(K_\Gamma)$ be its
$n$th coordinate matrix. Then $A$ is invertible over $\BB_\Gamma$ if and
only if all $A_n$ are invertible and
\[
 \{v(\det A_n):n\ge1\}
\]
is finite. Equivalently, for a matrix $A\in M_d(\HH_\Gamma)$ there exist
$B,C\in M_d(\HH_\Gamma)$ with $AB+hC=I_d$ exactly under these same
conditions at the $a_n$.
\end{corollary}
\begin{proof}
Over a commutative ring a square matrix is invertible exactly when its
determinant is a unit. Apply \cref{thm:units} to the determinant and use
the adjugate formula. The final assertion follows by lifting a matrix
inverse through \cref{thm:quotient}; a right inverse of a square matrix
over this commutative ring already forces a unit determinant.
\end{proof}
This matrix statement is a determinant consequence, not an independent
claim of a new matrix inversion theorem. Its useful content is the exact
support condition on infinitely many surcomplex fibres.

\subsection{The B\'ezout criterion and extension invariance}
\begin{proof}[Proof of \cref{thm:main-bezout}]
The identity $hU+FV=1$ exists precisely when the class of $F$ is a unit
in $\HH_\Gamma/(h)$. By \cref{thm:quotient}, this class corresponds to
$(F(a_n))_n$. Apply \cref{thm:units}.

Now let $\Gamma\subset\Gamma'$ be an ordered extension. Evaluation of
$F$ at each ordinary $a_n$ does not change. Its nonzero valuations are
the same elements, with the same order and the same cardinality, in the
larger group. The criterion is therefore unchanged. In particular, a
solution that exists in the larger coefficient ring already exists in
the original one.
\end{proof}

Here is also a direct construction in the positive case. Use
\eqref{eq:inverse-support} to build the inverse value family, lift it to
$V$ by interpolation, and then set
\[
 U=\frac{1-FV}{h}.
\]
The numerator vanishes coefficientwise at all $a_n$, so the quotient is
holomorphic coefficientwise on the same $D$. This construction is
noncanonical because ordinary interpolation is noncanonical, but its
support admissibility is proved, not presumed.

\begin{remark}[Why a minimum over several generators is insufficient]
The two-generator theorem treats the pair $(h,F)$ with the fixed ordinary
divisor $h$. It is not a criterion for an arbitrary row of elements in
$\BB_\Gamma$ obtained by replacing $v(x_n)$ with a pointwise minimum.
For example, if $x_n=t^{\gamma_n}$, then the two elements $x$ and
$x+t^3$ generate the unit ideal, since their difference is the unit $t^3$,
even though their pointwise minimum valuations form an infinite set.
\end{remark}

\section{The explicit uniform corona counterexample}\label{sec:corona}

\subsection{Admissibility and the unavoidable descending support}
Since $\gamma_n=2-1/n$ is strictly increasing, its image is well ordered
of order type $\omega$. Therefore \eqref{eq:F-main} is an element of
$\HH$, with $\vH(F)=1$ and $F(a_n)=t^{\gamma_n}$. The valuation set is
infinite, so \cref{thm:main-bezout} gives
\begin{equation}\label{eq:proper-I}
 I=(h,F)\subsetneq\HH.
\end{equation}
Equivalently, any putative identity evaluated at $a_n$ would force
\[
 V(a_n)=t^{-\gamma_n}=t^{-2+1/n}\qquad(n\ge1).
\]
Since an ordinary evaluation only deletes support exponents, all
$-2+1/n$ would have to belong to $\supp V$. This is impossible. The
argument remains unchanged in any larger ordered group, including
set-sized groups of surreal exponents of arbitrarily high rank.

\subsection{All actual surcomplex points satisfy the lower bound}
\begin{lemma}[Simple ordinary zeros do not acquire extra monad zeros]
\label{lem:hzeros}
For every ordered group extension containing the coefficients,
$h(c+\varepsilon)=0$ with $c\in D$ and $v(\varepsilon)>0$ holds exactly
when $c=a_n$ and $\varepsilon=0$ for some $n$.
\end{lemma}
\begin{proof}
If $h(c)\ne0$, its nonzero constant Taylor term prevents vanishing.
At a simple zero $a_n$,
\[
 h(a_n+\varepsilon)=h'(a_n)\varepsilon(1+\eta),
 \qquad v(\eta)>0,
\]
when $\varepsilon\ne0$. The last factor is a unit, and the others are
nonzero. The case $\varepsilon=0$ is immediate.
\end{proof}
Thus $h$ and $F$ have no common zero at any actual surcomplex point over
$D$, not merely no common zero at the ordinary centers.

The uniform norm estimate is stronger. Differentiating coefficientwise
does not lower Hahn exponents, so Taylor expansion gives, for
$\varepsilon\ne0$ and $q=v(\varepsilon)>0$,
\begin{equation}\label{eq:F-lipschitz}
 v\bigl(F(c+\varepsilon)-F(c)\bigr)\ge1+q.
\end{equation}
For clarity, this uses only that every positive-order Taylor term of an
ordinary coefficient contains a factor $\varepsilon$, while all
exponents of $F$ are at least $1$.

\begin{proposition}[The sharp lower bound]\label{prop:sharp-corona}
For all $z\in D^\sharp_\R$,
\begin{equation}\label{eq:strictlower}
 \max\{|h(z)|_v,|F(z)|_v\}>e^{-2}.
\end{equation}
The infimum of the left side is exactly $e^{-2}$ and is not attained.
\end{proposition}
\begin{proof}
Write $z=c+\varepsilon$. If $c\notin\{a_n\}$, then $v(h(z))=0$, so
the maximum is at least $1$. If $c=a_n$ and $\varepsilon=0$, then
$h(z)=0$ and $v(F(z))=\gamma_n<2$.

It remains to consider $c=a_n$ and $\varepsilon\ne0$. Put
$q=v(\varepsilon)$. Simplicity gives $v(h(z))=q$. If $q<2$, this already
proves \eqref{eq:strictlower}. If $q\ge2$, then
\eqref{eq:F-lipschitz} gives
\[
 v\bigl(F(z)-t^{\gamma_n}\bigr)\ge1+q\ge3>\gamma_n,
\]
so $v(F(z))=\gamma_n<2$. This proves the strict lower bound in all cases.
At $z=a_n$, the maximum is $e^{-\gamma_n}$, which tends to $e^{-2}$.
Thus the stated infimum is exact and is never attained.
\end{proof}

By \cref{prop:banach}, $\norm h_H=1$ and $\norm F_H=e^{-1}$.
By \eqref{eq:proper-I} and \cref{lem:distance}, the distance of $1$ from
$I$ is one. Hence even \emph{approximate} B\'ezout identities cannot have
uniform residual norm smaller than one.

This does not contradict Carleson's corona theorem
\cite{Carleson}. The scalar field, absolute value, and function algebra
are different. Our coefficients are not required to be bounded in the
ordinary complex norm, and the norm here treats all nonzero complex
constants as units of size one. Nor does it contradict results for
convergent Hahn-meromorphic expansions such as \cite{MullerStrohmaier},
whose function class is different.

\subsection{Local solutions and failure of principality}
\begin{proposition}[Local B\'ezout solutions on ordinary neighborhoods]
\label{prop:local-bezout}
Every $c\in D$ has an ordinary neighborhood $W\subset D$ on which there
are $U_W,V_W\in\OO(W)((t^\R))$ satisfying $hU_W+FV_W=1$.
\end{proposition}
\begin{proof}
If $c$ is not a zero of $h$, shrink $W$ so $h$ is nowhere zero and take
$U_W=1/h$, $V_W=0$. At $c=a_n$, choose $W$ containing no other zero of
$h$. Put
\[
 V_W=t^{-\gamma_n},\qquad
 U_W=\frac{1-t^{-\gamma_n}F}{h}.
\]
Each coefficient of the numerator vanishes at $a_n$, and division by the
single simple zero gives holomorphic coefficients on $W$. Translation
of the support by $-\gamma_n$, together with a possible exponent zero,
preserves well ordering. Thus these are legitimate local Hahn sections.
\end{proof}

\begin{proposition}[A proper, nonprincipal, two-generated free ideal]
\label{prop:nonprincipal}
The ideal $I=(h,F)$ is not principal.
\end{proposition}
\begin{proof}
Suppose $I=(g)$. Then $h=gu$ and $F=gw$ for some $u,w\in\HH$. Let
$\alpha=\vH(g)$. Since $\vH(h)=0$, one has $\vH(u)=-\alpha$, and the
leading coefficients satisfy
\[
 g_\alpha u_{-\alpha}=h.
\]
Every zero of $g_\alpha$ lies among the $a_n$. If $g_\alpha(a_n)=0$,
then simplicity of the zero of $h$ forces $u_{-\alpha}(a_n)\ne0$.
Consequently $u(a_n)\ne0$. Evaluating $h=gu$ at $a_n$ gives $g(a_n)=0$,
and evaluating $F=gw$ then gives $F(a_n)=0$, a contradiction. Hence
$g_\alpha$ has no zeros on $D$. By \cref{lem:global-unit}, $g$ is a unit,
contradicting \eqref{eq:proper-I}.
\end{proof}

The term \emph{free ideal} here means an ideal whose elements have no
common geometric zero. In classical ordinary holomorphic function rings,
free ideals and their maximal extensions are old objects, but finitely
generated ideals are B\'ezout and a proper finitely generated free ideal
cannot occur \cite{Henriksen52,Henriksen53}. Our explicit ideal is
therefore a genuine change in finite ideal theory caused by Hahn support,
not merely a restatement that classical global rings have free maximal
ideals.

\subsection{Every finite observation misses the obstruction}
For a finite set $J\subset\N$, the inverse prescriptions
$V(a_n)=t^{-\gamma_n}$ for $n\in J$ have finite common support. They are
realized, for example, by
\[
 V_J(z)=\sum_{n\in J}t^{-\gamma_n}e_n(z).
\]
Thus the obstruction cannot be detected by checking any finite subset of
the center values. Increasing the number of checked centers forces a
strictly decreasing sequence of new inverse exponents. Finite exact
computation can exhibit arbitrarily long initial portions of that
sequence, but its incompatibility with well ordering is an infinite
mathematical argument.

Combining the preceding propositions proves
\cref{thm:main-corona} in full.

\section{Growth gauges and prime ideals}\label{sec:prime-gauges}
We now turn to a different manifestation of global compatibility. Fix
$\Gamma=\R$ and a nonprincipal ultrafilter $\UU$ on $\N_{\ge1}$.
Existence of this ultrafilter is the only specifically ultrafilter-based
choice input. The explicit corona counterexample does not use it.

\subsection{Why a uniform lower valuation bound matters}
For every $G\in\HH$, \eqref{eq:eval-bound} gives
\begin{equation}\label{eq:uniformlower}
 v(G(a_n))\ge\vH(G)\quad\text{for all }n.
\end{equation}
This is a bound for one fixed function, not a bound independent of $G$.
It prevents multiplication by an arbitrary element of $\HH$ from
cancelling an arbitrarily fast growth rate of valuations.

\begin{definition}[Growth-gauge ideal]\label{def:gauge}
For a function $w:\N_{\ge1}\to\R_{\ge1}$, define
\begin{equation}\label{eq:gauge-ideal}
 P_w=\{G\in\HH:\{n:v(G(a_n))>m w(n)\}\in\UU
                         \text{ for every integer }m\ge1\}.
\end{equation}
In words, the valuations grow faster than every constant multiple of
$w$ along the chosen ultrafilter.
\end{definition}

\begin{theorem}[Prime ideals from arbitrary positive gauges]
\label{thm:gauge-prime}
For every $w:\N\to\R_{\ge1}$, the set $P_w$ is a proper prime ideal of
$\HH$ containing $h$. If $w'(n)/w(n)$ tends to $+\infty$ along $\UU$,
in the sense that
\[
 \{n:w'(n)>M w(n)\}\in\UU\qquad\text{for every }M>0,
\]
then $P_{w'}\subsetneq P_w$.
\end{theorem}
\begin{proof}
The zero function belongs to $P_w$, and additive inverses do not change
valuations. If $G,H\in P_w$, intersect the two ultrafilter-large sets on
which their valuations exceed $m w(n)$. On that intersection the
ultrametric inequality gives the same bound for $G+H$.

Let $A\in\HH$ be nonzero and $G\in P_w$. Put $L=\vH(A)$ and choose an
integer $C\ge\max\{0,-L\}$. For every $m\ge1$, on an ultrafilter-large
set,
\[
 v(G(a_n))>(m+C)w(n),\qquad v(A(a_n))\ge L.
\]
Because $w(n)\ge1$, these imply $v((AG)(a_n))>m w(n)$, with the same
conclusion when a value is zero. Thus $P_w$ is an ideal. It is proper
because $v(1)=0$, and it contains $h$ because $h(a_n)=0$ for every $n$.

For primality use the contrapositive. If $G\notin P_w$ and $H\notin P_w$,
there are positive integers $m,k$ for which
\[
 \{n:v(G(a_n))\le m w(n)\}\in\UU,\qquad
 \{n:v(H(a_n))\le k w(n)\}\in\UU.
\]
The ultrafilter property supplies these complements. On their
intersection both values are nonzero and
\[
 v((GH)(a_n))\le(m+k)w(n).
\]
Hence $GH\notin P_w$, proving primality.

If $w'/w\to_\UU+\infty$, then $w'>w$ on an ultrafilter-large set, so
$P_{w'}\subseteq P_w$. To prove strictness prescribe
\[
 W(a_n)=t^{\lceil w'(n)\rceil}.
\]
All required exponents lie in $\N_{\ge1}$, a well-ordered set, so
\cref{thm:quotient} supplies $W\in\HH$. Since
$\lceil w'(n)\rceil\ge w'(n)$, it belongs to $P_w$. Since
$\lceil w'(n)\rceil\le2w'(n)$, it does not belong to $P_{w'}$.
\end{proof}

\begin{remark}[Dependence on growth class]
If positive constants $c,C$ satisfy $cw(n)\le w'(n)\le Cw(n)$ on an
ultrafilter-large set, then $P_w=P_{w'}$. This follows by adjusting the
integer multiplier $m$ in the definition. Thus the construction detects
asymptotic growth classes, not a chosen numerical normalization of the
gauge.
\end{remark}

The use of growth rates to construct prime chains has strong precedents.
Henriksen associated prime ideals of the ring of entire functions with
growth of ordinary zero multiplicities \cite{Henriksen53}; later
power-series constructions are discussed in \cite{KangToan}. Here the
growing quantity is instead the \emph{Hahn valuation of a nonzero value}
at a fixed simple divisor. The next section identifies the additional
relative information this supplies.

\section{An infinite-dimensional spectral fibre}\label{sec:fibres}

\subsection{A continuum chain}
For $r\ge0$ put $w_r(n)=n^r$ and write $P_r=P_{w_r}$. For $r<s$,
$n^{s-r}\to+\infty$ in the ordinary sense, hence along every
nonprincipal ultrafilter. By \cref{thm:gauge-prime},
\begin{equation}\label{eq:prime-chain}
 P_s\subsetneq P_r\qquad(0\le r<s).
\end{equation}
Explicit separating functions $W_s$ may be chosen with
\begin{equation}\label{eq:prime-witness}
 W_s(a_n)=t^{\lceil n^s\rceil},\qquad
 W_s\in P_r\setminus P_s\quad(r<s).
\end{equation}
For $s>0$, one especially concrete representative is
\[
 W_s(z)=\sum_{k\ge1}t^k
          \sum_{\{n:\lceil n^s\rceil=k\}}e_n(z).
\]
Each inner sum is finite. For $s=0$ simply take $W_0=t$.

\subsection{The contraction is one classical maximal ideal}
\begin{proposition}\label{prop:contraction}
For every gauge $w\ge1$,
\begin{equation}\label{eq:contraction}
 P_w\cap\OO(D)=\mm_\UU
   =\{f\in\OO(D):\{n:f(a_n)=0\}\in\UU\}.
\end{equation}
The ideal $\mm_\UU$ is maximal and is not evaluation at any ordinary
point of $D$. Its residue field is naturally $\C^\N/\UU$.
\end{proposition}
\begin{proof}
For an ordinary holomorphic $f$, each $f(a_n)$ is either zero, with
valuation $+\infty$, or a nonzero ordinary complex number, with valuation
zero. The condition in \eqref{eq:gauge-ideal} is therefore exactly the
condition defining $\mm_\UU$.

By \cref{lem:ordinary-interp}, ordinary evaluation maps $\OO(D)$ onto
$\C^\N$. Quotienting $\C^\N$ by equality on an ultrafilter-large set
produces a field: a nonzero class is nonzero on such a set, and its inverse
is obtained by taking reciprocal coordinates there and arbitrary values
elsewhere. Thus the composite map onto $\C^\N/\UU$ has kernel
$\mm_\UU$, proving maximality.

The ideal contains $h$, so any common ordinary zero would have to be an
$a_n$. But $e_n\in\mm_\UU$ since it vanishes at all $a_m$ except $a_n$,
while $e_n(a_n)=1$. Thus it has no common ordinary zero.
\end{proof}

\subsection{The ordinary extended ideal and an explicit maximal ideal}
There are two different ways to pass from the classical maximal ideal
$\mm_\UU$ to Hahn families. The first is ordinary ideal extension,
which uses finite sums. The second is coefficientwise membership,
which allows all coefficients to lie in the ideal independently. They
are not the same operation.

\begin{theorem}[Two ideals over the same classical point]
\label{thm:two-ideals}
Put $k_\UU=\OO(D)/\mm_\UU\simeq\C^\N/\UU$. Define
\[
 Q_\UU=\mm_\UU\HH,\qquad
 J_\UU=\left\{\sum_\gamma f_\gamma t^\gamma:
                           f_\gamma\in\mm_\UU\text{ for all }\gamma\right\}.
\]
Then:
\begin{enumerate}
\item $Q_\UU=\{G:\{n:G(a_n)=0\}\in\UU\}$ is prime. In particular,
      the spectral fibre $\HH/\mm_\UU\HH$ is a domain and embeds into
      the field $K^\N/\UU$.
\item Coefficientwise reduction induces a surjection
      $\rho:\HH\twoheadrightarrow k_\UU((t^\R))$ with kernel $J_\UU$.
      Thus $J_\UU$ is an explicit maximal ideal.
\item For every finite real $r\ge0$,
      $Q_\UU\subsetneq P_r\subsetneq J_\UU$.
      The corona ideal $(h,F)$ is contained in $J_\UU$.
\end{enumerate}
\end{theorem}
\begin{proof}
If $G=\sum_{j=1}^m f_jG_j$ with $f_j\in\mm_\UU$, the common zero set
of the finitely many $f_j(a_n)$ belongs to $\UU$. Therefore
$G(a_n)=0$ on an ultrafilter-large set.

Conversely, suppose $G(a_n)=0$ for $n$ in some $E\in\UU$.
By ordinary interpolation choose $f\in\OO(D)$ with $f(a_n)=0$ on $E$
and $f(a_n)=1$ off $E$. Then $f\in\mm_\UU$, and $(1-f)G$ vanishes at
every $a_n$. By \cref{thm:quotient}, $(1-f)G=hB$ for some $B\in\HH$.
Since both $f$ and $h$ belong to $\mm_\UU$, the identity
$G=fG+hB$ proves $G\in Q_\UU$. The characterization is exactly the
kernel of evaluation into the ultraproduct field $K^\N/\UU$, proving
primality and the embedding assertion.

Reducing each coefficient modulo $\mm_\UU$ is a homomorphism by the
finite coefficient sums in Hahn multiplication. Every Hahn series over
$k_\UU$ has a lift: choose a representative in $\OO(D)$ for each of its
nonzero coefficients, on the same well-ordered support. The target is
a Hahn field. Its kernel is exactly $J_\UU$, proving maximality.

An element of $Q_\UU$ has zero values on an ultrafilter-large set, so
belongs to every $P_r$. This inclusion is strict because
$W_{r+1}\in P_r$ has nonzero values at every $a_n$, hence is not in
$Q_\UU$.

If $G=\sum f_\gamma t^\gamma\in P_r$, fix an exponent $\gamma$ and an
integer $m>\max\{0,\gamma\}$. On an ultrafilter-large set,
$v(G(a_n))>mn^r\ge m>\gamma$. Hence $f_\gamma(a_n)=0$ there. This shows
$f_\gamma\in\mm_\UU$ for every $\gamma$, so $P_r\subseteq J_\UU$.
The inclusion is strict because all $e_n$ lie in $\mm_\UU$, whence
$F\in J_\UU$, whereas $v(F(a_n))=\gamma_n<2$ gives $F\notin P_r$.
Finally $h\in\mm_\UU$ and $F\in J_\UU$, so $(h,F)\subseteq J_\UU$.
\end{proof}

\begin{theorem}[Infinite relative dimension]\label{thm:fibre}
The domain
\begin{equation}\label{eq:fibre-ring}
 \HH\otimes_{\OO(D)}k_\UU=\HH/\mm_\UU\HH
\end{equation}
has infinite Krull dimension and contains the continuum prime chain
induced by $(P_r)_{r\ge0}$. The algebra $\HH$ is not Noetherian, and
\[
 J_\UU\cap\OO(D)=\mm_\UU,\qquad \dim\HH_{J_\UU}=\infty.
\]
\end{theorem}
\begin{proof}
The identity \eqref{eq:fibre-ring} follows because $\mm_\UU$ is maximal.
The domain assertion is \cref{thm:two-ideals}. Each $P_r$ contains
$Q_\UU$ and contracts to $\mm_\UU$. The strict inclusions
\eqref{eq:prime-chain} survive in the quotient, giving arbitrarily long
finite prime chains. Their cardinality is $2^{\aleph_0}$.

For non-Noetherianity use the infinite ascending chain
\[
 P_1\subsetneq P_{1/2}\subsetneq P_{1/3}\subsetneq\cdots.
\]
This avoids the incorrect general inference that infinite Krull
dimension by itself forces a ring to be non-Noetherian. All the $P_r$
lie below the explicit maximal ideal $J_\UU$. Localization there
preserves strict inclusion, so $\dim\HH_{J_\UU}=\infty$. Its contraction
is $\mm_\UU$ directly from its coefficientwise definition.
\end{proof}

\begin{corollary}[Coefficientwise reduction is not scalar base change]
\label{cor:basechange}
The natural surjection
\[
 \HH\otimes_{\OO(D)}k_\UU\longrightarrow k_\UU((t^\R))
\]
has nonzero kernel $J_\UU/Q_\UU$. The class of $F$ is a specific nonzero
element of that kernel. Thus the Hahn-series construction does not
commute with this residue-field base change in the naive coefficientwise
sense.
\end{corollary}
\begin{proof}
Factor $\rho$ through $Q_\UU$. Its kernel is as stated. Every
coefficient of $F$ belongs to $\mm_\UU$, but $F$ has no zero value on
any $a_n$, so \cref{thm:two-ideals} gives $F\notin Q_\UU$.
\end{proof}

\subsection{All the growth primes have the same norm closure}
The topological behavior is as sharp as the algebraic behavior.

\begin{theorem}[Closure collapse of the prime chain]
\label{thm:closure}
In the complete norm $\norm{\cdot}_H$,
\[
 \overline{Q_\UU}=P_0,\qquad \overline{P_r}=P_0\quad(r\ge0).
\]
In particular $P_0$ is a closed nonmaximal prime ideal, while all $P_r$
with $r>0$ are nonclosed. Moreover
\begin{equation}\label{eq:distance-F}
 \operatorname{dist}(F,Q_\UU)
 =\operatorname{dist}(F,P_0)=e^{-2},
\end{equation}
and neither distance is attained.
\end{theorem}
\begin{proof}
Suppose $G\in\overline{Q_\UU}$. For every positive integer $m$ choose
$A\in Q_\UU$ with $\vH(G-A)>m$. On an ultrafilter-large set $A(a_n)=0$,
so $v(G(a_n))>m$ there. Thus $G\in P_0$.

Conversely, let $G=\sum f_\gamma t^\gamma\in P_0$. For a positive
integer $m$, let $G_{\le m}$ be its Hahn truncation to exponents at most
$m$. On the ultrafilter-large set where $v(G(a_n))>m$, all coefficients
of this truncation vanish at $a_n$. Hence $G_{\le m}\in Q_\UU$ by
\cref{thm:two-ideals}. The remainder is either zero or has valuation
strictly above $m$, so $G_{\le m}\to G$. This proves
$\overline{Q_\UU}=P_0$. Since $Q_\UU\subset P_r\subset P_0$,
the closure of every $P_r$ equals $P_0$. The strictness of the chain and
$P_0\subsetneq J_\UU$ give the closedness and nonmaximality statements.

Each finite partial sum $F_N=\sum_{n=1}^N t^{\gamma_n}e_n$ lies in
$Q_\UU$, since its coefficients lie in $\mm_\UU$ and the sum is finite.
The norm $\norm{F-F_N}_H=e^{-\gamma_{N+1}}$ tends to $e^{-2}$, proving
the upper bounds in \eqref{eq:distance-F}. For any $A\in P_0$, there is
an ultrafilter-large, hence nonempty, set on which $v(A(a_n))>2$.
At any such $n$,
\[
 v\bigl(F(a_n)-A(a_n)\bigr)=\gamma_n<2,
\]
so $\norm{F-A}_H\ge e^{-\gamma_n}>e^{-2}$. This gives the lower bounds
and nonattainment for $P_0$, and therefore for its subset $Q_\UU$.
\end{proof}

The truncation $G_{\le m}$ may contain infinitely many exponents; it is
not being approximated by finitely many coefficients. This point is
essential when a well-ordered support has a finite accumulation point.
The finite partial sums $F_N$ therefore do not contradict
\cref{thm:closure}: their limit distance from $F$ is positive.

The assertion is stronger than the statement that a global holomorphic
ring has large dimension. Classical holomorphic rings already have free
prime chains \cite{Henriksen53}. Our chain lies \emph{over one and the
same classical maximal ideal}; its nontrivial order is created in the
Hahn extension, not inherited from a chain in $\OO(D)$.

\subsection{The chain is invisible at all geometric points}
\begin{proposition}\label{prop:invisible}
For each $r\ge0$, no $a\in D^\sharp_\R$ satisfies
$P_r\subseteq\ker\ev_a$. The same remains true for evaluation in any
larger surcomplex workspace.
\end{proposition}
\begin{proof}
Every $P_r$ contains $h$. Thus such an $a$ would have to be one of the
exact points $a_n$, by \cref{lem:hzeros}. Every $P_r$ also contains every
$e_n$: its values are zero on a cofinite, hence ultrafilter-large, set.
But $e_n(a_n)=1$. This excludes $a_n$ and completes the proof.
\end{proof}

Consequently the explicit maximal ideal $J_\UU$, and every other maximal
ideal containing $P_0$, is genuinely nongeometric, not evaluation at a
point with an overlooked infinitesimal displacement. The corona ideal $I$ and the ideals $P_r$ illustrate
\emph{different} mechanisms: the values of $F$ have bounded valuations,
so $F\notin P_r$ for every $r\ge0$. The displayed growth primes do not contain the corona ideal, although
the common maximal ideal $J_\UU$ contains both that ideal and the
entire chain.

\section{Every geometric localization is a discrete valuation ring}
\label{sec:dvr}
We return to a general ordered abelian group $\Gamma$. This section
proves the regular local behavior rather than inferring it from a global
Noetherian statement. Denote the ordinary coordinate function by $z$ and
coefficientwise differentiation by $\partial_z$.

\subsection{Global division by an infinitesimally shifted coordinate}
Fix $a=c+\varepsilon\in D^\sharp_\Gamma$. For $f\in\OO(D)$, set
$b_j=f^{(j)}(c)/j!$ and define
\begin{equation}\label{eq:remainders}
 R_kf(z)=\frac{f(z)-\sum_{j=0}^k b_j(z-c)^j}{(z-c)^{k+1}},
 \qquad k\ge0.
\end{equation}
The singularity at $c$ is removable, so $R_kf\in\OO(D)$.

\begin{lemma}[Support-controlled divided difference]\label{lem:division}
For every $F\in\HH_\Gamma$ there is $Q\in\HH_\Gamma$ such that
\begin{equation}\label{eq:division}
 F-F(a)=(z-a)Q.
\end{equation}
If $\varepsilon\ne0$, $S=\supp F$, and $E=\supp\varepsilon$, one may
take
\begin{equation}\label{eq:division-explicit}
 Q(z)=\sum_{\gamma\in S}\sum_{k\ge0}
          t^\gamma\varepsilon^k R_kf_\gamma(z),
 \qquad \supp Q\subseteq S+E^*.
\end{equation}
For $\varepsilon=0$, use only $k=0$.
\end{lemma}
\begin{proof}
The formula is strongly summable by \cref{lem:neumann}. At each Hahn
exponent its coefficient is a finite sum of functions holomorphic on the
same $D$, so $Q\in\HH_\Gamma$.

To verify the identity, first treat one coefficient $f$. The remainders
satisfy
\[
 (z-c)R_0f=f-b_0,\qquad
 (z-c)R_kf-R_{k-1}f=-b_k\quad(k\ge1).
\]
Multiplying $\sum_{k\ge0}\varepsilon^kR_kf$ by
$z-c-\varepsilon$ and collecting the finitely justified coefficient
contributions gives
\[
 f(z)-\sum_{k\ge0}b_k\varepsilon^k=f(z)-f(a).
\]
Multiply by $t^\gamma$ and sum over the support of $F$.
\end{proof}

\begin{corollary}[Evaluation ideals]\label{cor:eval-ideal}
For every $a\in D^\sharp_\Gamma$,
\[
 \ker\ev_a=(z-a),\qquad \HH_\Gamma/(z-a)\simeq K_\Gamma.
\]
Thus $\mm_a=(z-a)$ is a maximal ideal.
\end{corollary}
\begin{proof}
Evaluation is surjective by \cref{prop:eval}, and
\cref{lem:division} proves the kernel formula.
\end{proof}

\subsection{A finite order of vanishing at every geometric point}
\begin{lemma}\label{lem:finite-order}
If $F\in\HH_\Gamma$ is nonzero and $a=c+\varepsilon$, then
$\partial_z^jF(a)\ne0$ for some finite $j$. If the leading ordinary
coefficient $f_\alpha$ has multiplicity $m$ at $c$, then one may take
$j=m$.
\end{lemma}
\begin{proof}
The ordinary holomorphic function $f_\alpha$ is nonzero, so its order of
vanishing at $c$ is finite. The coefficient at exponent $\alpha$ of
$\partial_z^mF(a)$ is $f_\alpha^{(m)}(c)\ne0$. Positive powers of
$\varepsilon$ have strictly positive exponents, and all higher Hahn
coefficients start above $\alpha$, so they cannot cancel this term.
\end{proof}

\begin{theorem}[Geometric discrete valuation rings]\label{thm:dvr}
For $a\in D^\sharp_\Gamma$, let $d_a(F)$ be the least $j\ge0$ such that
$\partial_z^jF(a)\ne0$. Every nonzero $F\in\HH_\Gamma$ factors uniquely
in the form
\begin{equation}\label{eq:localfactor}
 F=(z-a)^{d_a(F)}G,\qquad G(a)\ne0,
\end{equation}
with respect to its exponent $d_a(F)$. The localization
$(\HH_\Gamma)_{\mm_a}$ is a discrete valuation ring with uniformizer
$z-a$, residue field $K_\Gamma$, and Krull dimension one.
\end{theorem}
\begin{proof}
The least $j$ exists by \cref{lem:finite-order}. If $F(a)=0$, divide by
$z-a$ using \cref{lem:division}. The coefficientwise product rule gives
\[
 \partial_z^j((z-a)Q)(a)=j\,\partial_z^{j-1}Q(a)\quad(j\ge1).
\]
Thus the least nonzero derivative order decreases by one at each
step, and after $d_a(F)$ steps one obtains \eqref{eq:localfactor} with
$G(a)=\partial_z^{d_a(F)}F(a)/d_a(F)!\ne0$.

A fraction in the localization has denominator nonvanishing at $a$.
Its numerator has the displayed factorization, and $G$ becomes a unit.
Therefore every nonzero element of the local ring is a unit times a
unique nonnegative power of $z-a$. Every nonzero ideal has an element
of least such exponent and is generated by that power. The maximal
ideal is nonzero and generated by $z-a$, proving that this is a discrete
valuation ring, of dimension one. The residue field was identified in
\cref{cor:eval-ideal}.
\end{proof}

\begin{corollary}[Completion and finite jets]\label{cor:completion}
For every $N\ge1$, Taylor jets identify
\[
 (\HH_\Gamma)_{\mm_a}/(z-a)^N
       \simeq K_\Gamma[X]/(X^N).
\]
Consequently
\[
 \widehat{(\HH_\Gamma)_{\mm_a}}\simeq K_\Gamma[[X]].
\]
\end{corollary}
\begin{proof}
Repeated divided differences express any numerator modulo $(z-a)^N$ as
its finite Taylor polynomial with coefficients in $K_\Gamma$. A
nonvanishing denominator has a unit constant jet and is invertible in
the truncated polynomial ring. Conversely every polynomial jet is
represented by a polynomial in $z-a$ with coefficients in the embedded
constant field. The kernel is exactly $(z-a)^N$, by
\cref{thm:dvr}. Taking the inverse limit gives the completion.
\end{proof}

The discrete valuation $d_a$ is an order of vanishing in the spatial
variable. It is distinct from the Hahn valuation $v$ on scalar values
and from $\vH$ on coefficient families. Confusing these three valuations
would obscure the central contrast: local spatial dimension is one,
while global Hahn-value growth creates an infinite-dimensional spectral
fibre elsewhere.

Together, \cref{thm:fibre,prop:invisible,thm:dvr,cor:completion} prove
\cref{thm:main-spectrum}.

\section{Real forms and the role of the value group}\label{sec:real}

\subsection{A surreal-valued version}
Let
\[
 \OO_{\R}(D)=\{f\in\OO(D): f(\overline z)=\overline{f(z)}\},\qquad
 \HH_{\R,\Gamma}=\OO_{\R}(D)((t^\Gamma)).
\]
Its constant field is $K_{\R,\Gamma}=\R((t^\Gamma))$, which embeds in
$\No$ for $\Gamma\subset\No$. Evaluation is defined at
$c+\varepsilon$ with $c\in(-1,1)$ and real infinitesimal
$\varepsilon\in K_{\R,\Gamma}$.

\begin{corollary}[Real coefficient form]\label{cor:real}
The explicit pair $h,F$ lies in $\HH_{\R,\R}$ and has no common zero
on the surreal thickening of $(-1,1)$. It satisfies the same sharp
valuation-corona bound there, and has no B\'ezout identity even after
complexification. The real algebra also has the prime chains and
infinite-dimensional fibres constructed above, and its geometric
localizations at real points are discrete valuation rings with residue
field $\R((t^\R))$.
\end{corollary}
\begin{proof}
All $a_n$ are real, and $h,e_n$ are invariant under complex conjugation.
The bounds in \cref{prop:sharp-corona} remain valid on the real monads,
and their sharpness is witnessed by the real points $a_n$.
Nonexistence of a B\'ezout identity in the larger complex ring implies
nonexistence in the real subring.

For real interpolation values, symmetrize an ordinary interpolant by
replacing $f(z)$ with
$\bigl(f(z)+\overline{f(\overline z)}\bigr)/2$. The interpolation quotient
therefore holds with $\R^\N$ in place of $\C^\N$. The valuation and
ultrafilter arguments are unchanged, with residue field
$\R^\N/\UU$ at the classical free maximal ideal. Finally,
\eqref{eq:remainders} preserves the real form when $c$ and
$\varepsilon$ are real, so the discrete valuation ring proof carries
over verbatim.
\end{proof}

\subsection{Discrete and nondiscrete rank-one groups}
The example $\gamma_n=2-1/n$ uses rational exponents. Thus the algebraic
nonexistence and the bounded corona phenomenon already occur over
$\C((t^\Q))$, not only over $\C((t^\R))$. More generally any ordered
value group with an infinite increasing bounded sequence admits the
same reciprocal-support obstruction in $\BB_\Gamma$.

For a discrete rank-one group $\Gamma=\Z\delta$, $\delta>0$, a nonempty
well-ordered set of valuations is bounded below. If it is also bounded
above, it is finite. Hence, for the particular pair $(h,F)$ in
\cref{thm:main-bezout}, a uniform upper bound on $v(F(a_n))$ does imply
solvability when all these values are nonzero. The bounded-spectrum
counterexample is therefore genuinely sensitive to the value group.

This does \emph{not} restore unrestricted global B\'ezout theory for a
discrete group. The legitimate series
\[
 F_{\mathrm{unb}}=\sum_{n\ge1}t^{n\delta}e_n
\]
still has $F_{\mathrm{unb}}(a_n)=t^{n\delta}$ and gives a proper
zero-free two-generated ideal. Its inverse value norms, however, are
unbounded. Likewise the prime-chain witnesses use only nonnegative
integer exponents; after multiplying exponents by $\delta$, the
infinite-dimensional fibre construction already exists in the ordinary
Laurent-series algebra $\OO(D)((t^\Z))$.

These distinctions separate two effects. Infinitely many allowable
leading scales suffice for global ideal-theoretic failure. A bounded
infinite set of those scales is what turns that failure into the
uniform valuation-corona counterexample.

\subsection{No repair by a larger surreal universe}
Any particular proposed solution $U,V$ expressed by Hahn coefficient
families uses a set of surreal exponents. The union of that set with the
exponents of the input is contained in a set-sized ordered subgroup of
$\No$. Extension invariance in \cref{thm:main-bezout} then applies. Thus
allowing additional infinite or infinitesimal scales cannot repair the
identity as long as the solution still has one well-ordered global
Hahn support. A change of the function class, not merely of the size of
the workspace, would be required.

\section{Consequences for the repository's analytic geometry}
\label{sec:repository}

\subsection{An assertion that needs a fixed-domain correction}
The inspected repository snapshot is
\begin{center}\small\ttfamily
 aa846271b4dcae2c055b216126a87210292ec19b.
\end{center}
The source \texttt{docs/surcomplex/analytic-geometry/article.tex},
in its discussion of four coefficient rings near source lines
213--223, includes the \emph{fixed-polydisk} ring
$\OO(D)((t^\Gamma))$ in a structure description asserting
Noetherianity and a description of maximal ideals by infinitesimal
evaluation ideals \cite{RepoGeometry}. For that fixed-domain ring,
the description is false as written.

The failure is not merely a missing point at another ordinary center:
\cref{thm:fibre} gives non-Noetherianity and a localization of infinite
Krull dimension, while \cref{prop:invisible} excludes all finite
surcomplex points over the disk, including infinitesimally displaced
ones. \Cref{thm:main-corona} also refutes a geometric
Nullstellensatz for arbitrary finitely generated ideals of this global
ring. Indeed, the proper ideal $I=(h,F)$ has empty geometric zero set,
so its geometric vanishing ideal is the whole ring, whereas
$\sqrt I$ is proper.

\subsection{What is not refuted}
The same repository carefully distinguishes the fixed-domain algebra
from common-domain \emph{germs}, radius-free coefficient germs, and
formal coefficient rings. Our counterexample does not transfer to any
of those rings just because their notation looks similar. It uses
infinitely many ordinary centers approaching the boundary, a resource
that disappears when one localizes around a single ordinary center.

Nor does it refute the finite-deformation theorem for a leading ordinary
system with a \emph{finite isolated zero scheme}
\cite{RepoFinite}. Our function $h$ has infinitely many ordinary zeros.
The finite-isolated-zero hypothesis is therefore absent. At each
individual geometric point, the positive regularity statement is
confirmed by our independent proof of \cref{thm:dvr}.

The global-divisor report already emphasizes support-restricted
interpolation and explicitly refrains from calling the fixed-disk ring
a PID \cite{RepoGlobal}. The present article should be read as a
continuation of that global-support viewpoint and a correction to the
more expansive fixed-ring attribution elsewhere, not as a wholesale
rejection of the repository's analytic framework.

\subsection{A replacement statement}
For the ring used here, a valid replacement is:
\begin{quote}
The fixed-domain Hahn algebra has geometric evaluation maximal ideals
with discrete valuation localizations and formal power-series
completions. These do not exhaust the global algebraic spectrum.
The algebra is non-Noetherian, has infinite-dimensional fibres over
classical free maximal ideals, and admits proper two-generated ideals
with no geometric zero, even under a uniform valuation lower bound.
\end{quote}
We have not proved that the global ring is or is not Jacobson, and no
claim about that separate property is implicit in the replacement.
Similarly, we do not claim to classify all abstract maximal ideals.

\section{Prior work, novelty boundaries, and further questions}
\label{sec:prior}

\subsection{What is classical, what is inherited, and what is proposed}
Hahn's construction and Neumann's support calculus provide the algebraic
summation rules \cite{Hahn,Neumann}. The inversion of $1+U$ for positive
support is classical; Poonen's 2026 note gives a new proof of that
classical result, not the common-support product criterion proved here
\cite{Poonen}. We use it only for the elementary geometric inverse.
The general formal-summability viewpoint is also established
\cite{BagayokoEtAl}.

Ordinary discrete interpolation, ordinary free maximal ideals, and
prime chains measured by growth have long histories
\cite{Henriksen52,Henriksen53,KangToan}. In particular, merely producing
an infinite-dimensional global function ring would not justify a
novelty claim. Our relative assertion fixes one classical maximal ideal
and constructs a continuum chain \emph{inside its fibre} by Hahn
valuations of nonzero values. We exhibit all the primes and all the
strictness witnesses explicitly once an ultrafilter is fixed.

The repository already contains global support obstructions and a more
general deformed interpolation quotient \cite{RepoGlobal}. The simple
quotient in \cref{thm:quotient}, ordinary interpolation, and the basic
idea that a descending support blocks gluing are consequently credited
as inputs or context. The proposed additional results are the exact
finite-leading-valuation criterion, the explicit \emph{normed} corona
counterexample with its sharp lower bound and unit distance, its
extension invariance, and the infinite-dimensional domain fibre
contrasted with all geometric localizations. The explicit maximal
ideal, failure of naive residue-field base change, and common norm
closure of the growth-prime chain sharpen that last contrast. The matrix criterion is included as
a useful consequence and is not elevated to a separate major discovery.

The literature check included primary sources on Hahn summability,
Hahn-meromorphic analysis, classical holomorphic ideal theory,
power-series prime chains, and the pinned repository documents. It did
not locate these precise fixed-domain results. This is a targeted
comparison, not a proof that no equivalent formulation exists in the
literature. No named open problem from the classical surreal literature
is asserted to have been solved.

\subsection{Natural next questions, not used as assumptions}
One can now ask for a support-theoretic criterion for a general finite
row $(F_1,\ldots,F_m)$ to be unimodular in $\HH$. The warning after
\cref{thm:main-bezout} shows why the minimum of their pointwise leading
valuations is not sufficient data. Finite cancellations among generators
must be recorded as well as leading scales.

A second question is whether the growth-gauge ideals exhaust a useful
specified subclass of the primes above $\mm_\UU$. They are not claimed
to exhaust the fibre. In particular, the bounded valuation obstruction
of the corona pair is not itself one of our faster-than-polynomial
conditions. Relating bounded support cuts to the remaining primes is a
separate problem.

A third question concerns replacement function algebras. The actual
geometric evaluations norm every individual element of $\HH$, yet they
do not detect whether a finite row generates the unit ideal. One may
seek a function class with stronger global support control or an
appropriately enlarged spectrum for which a corona principle is valid.
Such a replacement must specify its algebra and its topology; enlarging
the scalar field alone is ruled out by the extension theorem.

\subsection{Conclusion}
The central obstruction is exact: inverses at infinitely many centers
can demand a reverse well ordering incompatible with a single global
Hahn family. This yields a quantitative corona failure in a complete
uniform algebra, not merely a badly behaved formal prescription.
Independently, valuation growth produces long prime chains over one
classical free maximal ideal, while every actual surcomplex point still
has a regular one-dimensional local ring. A global surcomplex analytic
geometry based on these algebras must therefore keep its geometric
points, its ordinary-base germs, and its full algebraic spectrum
separate.

\appendix
\section{Proof dependencies and the scope of verification}
\label{app:audit}
The proof has a short dependency structure. The Hahn support lemma
justifies evaluation and positive-support inversion. Ordinary cardinal
interpolation gives the quotient by $h$. Exact units in that quotient
give the B\'ezout criterion and the bounded-spectrum counterexample.
The valuation inequality and an ultrafilter give the prime ideals;
integer-exponent interpolants prove strictness. A separate global
Taylor divided-difference identity gives every geometric discrete
valuation ring. No conclusion about global Noetherianity is used in
any local argument, and no local preparation theorem is needed in the
prime construction.

The archive includes \texttt{verify.py} and its recorded JSON output.
The script performs exact rational or polynomial checks of the
following finite components: the exponent differences and inverse
ordering; cardinal derivative values at finitely many centers;
polynomial divided-difference identities; the cases of the valuation
lower-bound argument on a rational grid; and finite common-support
inverse examples. It rejects Python's assertion-disabling mode.

These computations are consistency checks, not a formal verification of
the infinite theorems. In particular they do not prove the Hahn support
lemma, an ultrafilter property, interpolation of an arbitrary infinite
family, completeness, or any statement of historical priority. All
infinite arguments are supplied in the article. No Lean compilation or
proof-assistant certification is claimed.

\section{Notation at a glance}
\begin{longtable}{@{}L{0.25\textwidth}L{0.69\textwidth}@{}}
\toprule
Symbol & Meaning\\
\midrule
\endhead
$D$ & The fixed ordinary complex unit disk.\\
$\Gamma$ & A set-sized ordered abelian value group.\\
$K_\Gamma$ & The scalar Hahn field $\C((t^\Gamma))$.\\
$\HH_\Gamma$ & The fixed-domain coefficient algebra $\OO(D)((t^\Gamma))$.\\
$D^\sharp_\Gamma$ & Points $c+\varepsilon$ with $c\in D$ and
$\varepsilon$ infinitesimal in $K_\Gamma$.\\
$v$ & The minimum-support valuation of a scalar Hahn series.\\
$\vH$ & The minimum-support valuation of a Hahn family with holomorphic
coefficients.\\
$d_a$ & The finite spatial order of vanishing at a geometric point $a$;
not the Hahn valuation.\\
$\norm{\cdot}_H$ & For $\Gamma=\R$, the norm $e^{-\vH(\cdot)}$, equal to
the supremum valuation norm on $D^\sharp_\R$.\\
$\BB_\Gamma$ & Sequences of scalar Hahn values with one common
well-ordered support.\\
$h,a_n,e_n$ & The explicit simple divisor, its zeros, and its cardinal
functions in \eqref{eq:cardinal-data}.\\
$\gamma_n,F$ & The increasing bounded exponents $2-1/n$ and the Hahn
interpolant in \eqref{eq:F-main}.\\
$\UU$ & A fixed nonprincipal ultrafilter on the positive integers.\\
$P_w$ & The prime ideal of values growing faster in valuation than every
constant multiple of the gauge $w$.\\
$P_r$ & $P_w$ for $w(n)=n^r$.\\
$\mm_\UU$ & The classical free maximal ideal contracted from every
$P_r$.\\
$Q_\UU$ & The prime ideal $\mm_\UU\HH$, characterized by zero values
on an ultrafilter-large set.\\
$J_\UU$ & The maximal ideal of families with every coefficient in
$\mm_\UU$; kernel of reduction to $k_\UU((t^\R))$.\\
$\mm_a$ & The geometric evaluation maximal ideal $(z-a)$.\\
\bottomrule
\end{longtable}

\begin{thebibliography}{99}
\addcontentsline{toc}{section}{References}

\bibitem{Hahn}
H.~Hahn, \emph{\"Uber die nichtarchimedischen Gr\"o\ss ensysteme},
Sitzungsberichte der Kaiserlichen Akademie der Wissenschaften,
Mathematisch-Naturwissenschaftliche Klasse \textbf{116} (1907),
601--655.

\bibitem{Neumann}
B.~H.~Neumann, \emph{On ordered division rings},
Transactions of the American Mathematical Society \textbf{66} (1949),
202--252. \href{https://doi.org/10.1090/S0002-9947-1949-0032593-5}{doi:10.1090/S0002-9947-1949-0032593-5}.

\bibitem{Poonen}
B.~Poonen, \emph{Units in Hahn--Mal'cev--Neumann rings},
author preprint, January 14, 2026, 4 pp.
\url{https://math.mit.edu/~poonen/papers/malcev.pdf}.

\bibitem{BagayokoEtAl}
V.~Bagayoko, L.~S.~Krapp, S.~Kuhlmann, D.~Panazzolo, and M.~Serra,
\emph{Automorphisms and derivations on algebras endowed with formal
infinite sums}, arXiv:2403.05827v2 (2025).
\url{https://arxiv.org/abs/2403.05827}.

\bibitem{Gonshor}
H.~Gonshor, \emph{An Introduction to the Theory of Surreal Numbers},
London Mathematical Society Lecture Note Series 110,
Cambridge University Press, 1986.

\bibitem{vdDE}
L.~van den Dries and P.~Ehrlich, \emph{Fields of surreal numbers and
exponentiation}, Fundamenta Mathematicae \textbf{167} (2001), 173--188.
\href{https://doi.org/10.4064/fm167-2-3}{doi:10.4064/fm167-2-3}.

\bibitem{Henriksen52}
M.~Henriksen, \emph{On the ideal structure of the ring of entire
functions}, Pacific Journal of Mathematics \textbf{2} (1952), 179--184.
\url{https://msp.org/pjm/1952/2-2/pjm-v2-n2-p05-p.pdf}.

\bibitem{Henriksen53}
M.~Henriksen, \emph{On the prime ideals of the ring of entire functions},
Pacific Journal of Mathematics \textbf{3} (1953), 711--720.
\url{https://msp.org/pjm/1953/3-4/pjm-v3-n4-p04-p.pdf}.
The discussion on p.~712 explicitly notes the disk analogue of the
ordinary interpolation and B\'ezout arguments.

\bibitem{KangToan}
B.~G.~Kang and P.~T.~Toan, \emph{How to construct huge chains of prime
ideals in power series rings}, in M.~Fontana, S.~Frisch, and S.~Glaz
(eds.), \emph{Commutative Algebra}, Springer, 2014, 225--238.
\href{https://doi.org/10.1007/978-1-4939-0925-4_13}{doi:10.1007/978-1-4939-0925-4\_13}.
Publisher abstract and bibliography consulted; the subscription-only
chapter is not a proof dependency.

\bibitem{Carleson}
L.~Carleson, \emph{Interpolations by bounded analytic functions and the
corona problem}, Annals of Mathematics (2) \textbf{76} (1962), 547--559.
\href{https://doi.org/10.2307/1970375}{doi:10.2307/1970375}.

\bibitem{MullerStrohmaier}
J.~M\"uller and A.~Strohmaier, \emph{The theory of Hahn-meromorphic
functions, a holomorphic Fredholm theorem, and its applications},
Analysis \& PDE \textbf{7} (2014), no.~3, 745--770.
\href{https://doi.org/10.2140/apde.2014.7.745}{doi:10.2140/apde.2014.7.745}.
Preprint: \url{https://arxiv.org/abs/1205.0236}.

\bibitem{RepoGeometry}
\emph{Infinitesimal Analytic Geometry over the Surcomplex Numbers},
research draft in V.~Reshetnikov's \texttt{Surreal} repository,
\texttt{docs/surcomplex/analytic-geometry/article.tex},
commit \texttt{aa846271b4dcae2c055b216126a87210292ec19b},
particularly the four-ring discussion near source lines 213--223.
\href{https://github.com/VladimirReshetnikov/Surreal/blob/aa846271b4dcae2c055b216126a87210292ec19b/docs/surcomplex/analytic-geometry/article.tex}{Pinned source}.

\bibitem{RepoGlobal}
\emph{Global Support Obstructions: Moving Divisors, Mittag--Leffler,
and a Picard Dichotomy}, research draft in the same
repository and commit, \texttt{docs/surcomplex/global-divisors/}.
The report's README and its account of the support-restricted
Chinese remainder theorem in \S10 were consulted.
\href{https://github.com/VladimirReshetnikov/Surreal/blob/aa846271b4dcae2c055b216126a87210292ec19b/docs/surcomplex/global-divisors/README.md}{Pinned report overview}.

\bibitem{RepoFinite}
\emph{Finite Hahn Deformations: Division, Conservation of Multiplicity,
and Residue Duality}, research draft in the same repository
and commit, \texttt{docs/surcomplex/finite-deformations/article.tex},
abstract and \S1 (finite isolated leading zero scheme).
\href{https://github.com/VladimirReshetnikov/Surreal/blob/aa846271b4dcae2c055b216126a87210292ec19b/docs/surcomplex/finite-deformations/article.tex}{Pinned source}.

\end{thebibliography}
\end{document}
